\cotefile{156-4}
\scanpage{10}
\mspage{9}

Ainsi, on est \unsure{arrivé|amené} à la structure suivante

\textcircled{1} $\unsure{\mathcal{F}|F}_{\unsure{el|\text{él}}}$ \unsure{ens|ens.} des ``figures élémentaires''

\textcircled{2} relation d'ordre sur $\mathcal{F}_{el}$, $\leq$ (inclusion des figures)

\textcircled{3} relation symétrique réflexive $R$ \deleted{\unsure{$F,G$|$F\,R\,G$}} ``$F$ et $G$ \uline{\unsure{compatibles}}''.\ \deleted{\illegible{1}} \deleted{\unsure{Parties|Partie}}

\margin{\uline{Configuration} (\unsure{élémentaire}) partie de $\mathcal{F}_{\text{él}}$ \unsure{formée|fermée} d'él. deux \unsure{à|:} deux \unsure{compatibles|comp.tibles}}

% slanted
\margin{deux configurations sont dites compatibles si leur réunion l'est.}

\textcircled{4} \deleted{Antifiltre sur $\mathcal{F}_{el}$, des parties} Parmi les configurations, \deleted{\unsure{descripta|description|decripta}} \inserted{donnée de l'ens.} \deleted{\unsure{loc. finies}} \deleted{\unsure{prop.}} \deleted{\illegible{1}} de celles qui sont ``loc. finies'', avec \deleted{Avec} propriétés d'antifiltre \unsure{;|,} $\mathrm{Conf}_{\mathrm{l.f.}}\unsure{\subset|\in}\mathcal{P}(\mathcal{F}_{el})$

a) Toute configuration réduite \unsure{à|:} une seule figure él. est loc. finie

% slanted
\margin{Cas important (cas ``\uline{compact}'') : les configurations loc. finies sont les configurations finies.}

b) Si $\Phi$ est une configuration, $\Phi$ loc. finie ssi la \deleted{\unsure{réduite|réduit}} partie \unsure{fermée|fermée.} (pour la rel. d'ordre $\leq$) l'est

c) Si $\Phi$ est une configuration, \deleted{$\Phi=\bigcup_{i}$} loc. finie, alors tte $\unsure{\Psi|\psi}\subset\Phi$ est loc. finie

d) Si $\unsure{\Phi|\Psi}$ est une configuration, $\Phi=\bigcup_{i}\Phi_i$ (réunion finie) si les $\Phi_i$ sont loc. \unsure{finis|finies}, alors $\Phi$ aussi.
\scanpage{11}
\mspage{10}

On \unsure{definit|définit} alors l'\unsure{ens|ens.} $\unsure{\mathcal{F}|F}$ des \uline{figures} \inserted{$F$}, comme l'ens des parties $\Phi={\unsure{\Sigma|\overline{Z}}}'_{F}$ de $\mathcal{F}_{\unsure{el|\text{él}}}$, qui sont a) des configurations

b) \unsure{fermées|fermees} pour la relation d'ordre

c) loc finies. Figures \uline{finies} : parties finies de $\mathcal{F}_{el}$ satisfaisant a) et b).

% slanted
\margin{Figures compatibles : qui \unsure{le} sont \unsure{comme} configurations. La \uline{réunion} \inserted{des \unsure{$E_{el}$|Ens}} supports sera alors la figure somme, et l'intersection sera aussi une figure -- la figure intersection \ldots}

Notons comme propriétés, en plus des propriétés a) \unsure{à|:} d) de $4^\circ$) ci-dessus, la relation entre \textcircled{2} et \textcircled{3}

\deleted{Axiome \unsure{Si|de}}

\uline{Fig}$_1$ Si $\deleted{\unsure{\Phi|F}}\unsure{X|F}_1,\deleted{F}X_2\leq\deleted{F}X$ sont des fig. él. alors $\{\deleted{F}X_1,\deleted{F}X_2\}\in R$ i.e. $\{\deleted{F}X_1,\deleted{F}X_2\}$ est une configuration.

[Donc l'ens des $\deleted{F}X'\leq\deleted{F}X$ est une configuration, et en vertu de $4^\circ$ a) b) c'est \unsure{$\overline{\mathrm{m}}$|mê} une \uline{figure} : la figure \unsure{associée|associé} \unsure{à|:} $F$\unsure{\ldots|\,--}\ ]

% slanted
\margin{On l'appelle la figure élémentaire \unsure{associé|associée} à $F$}
\scanpage{12}
\mspage{11}

% vertical rule to the left of the condition
\uline{Fig}$_{\unsure{2|n}}$ \inserted{Soit $\deleted{F}\unsure{X|F}\in\unsure{\mathcal{F}|F}_{\unsure{el|\text{él}}}$,} \unsure{Considérons}\vadjust{\vbox to 0pt{\hbox{\kern-0.7em\vrule width 0.35pt height 0pt depth 3\baselineskip}\vss}} l'ens ordonné des $\deleted{F}X'\leq\deleted{F}X$ est de \uline{dim comb. finie}, i.e. il existe un entier $n$, tel que tte partie finie $\deleted{\unsure{F}}\Phi$ \deleted{\illegible{1}} majorée par $X$, \inserted{et tot. ordonnée} soit de cardinal $\leq n$.

La \uline{\unsure{dimension|dimention}} \uline{combinatoire} d'un $X\in\mathcal{F}_{el}$ est le Sup des cardinaux des parties finies $\Phi$ \unsure{précédentes}, \unsure{diminués|diminué} de 1. Donc la dim combinatoire est $\geq 0$, et elle est nulle ssi $X$ est un \uline{\unsure{élément|élémt}} \uline{minimal} de $\mathcal{F}_{\text{él}}$.

Dans les exemples \unsure{types|type}, \deleted{\unsure{F|les}} les figures élémentaires minimales correspondent aux parties fermées de $X$ qui sont des \uline{variétés} \uline{\unsure{connexes}}.

Si $\Phi\subset\mathcal{F}_{\text{él}}$ est une figure, \deleted{\unsure{sa}} \uline{dim combinatoire} est sa dim comb en tant qu'\unsure{ens|ens.} ordonné, c'est \unsure{donc|ou} le sup des $\mathrm{dim\,comb}\,X$, pour $X\in\Phi$. Elle est nulle ssi \deleted{$\Phi$ est \unsure{formé|formée} de \unsure{va}} \inserted{$\Phi\neq\varnothing$ et} la relation
\scanpage{13}
\mspage{12.}

d'ordre sur $\Phi$ est \unsure{discrète|discrète,} i.e. $\Phi\neq\varnothing$ et \unsure{formé|formée} d'éléments minimaux de $\unsure{\mathcal{F}|F}$.

Ce sont donc les \deleted{\unsure{ens.|ens} configurations} parties de $\mathcal{F}_{\unsure{e\bar l|\overline{el}|\text{él}}\,\mathrm{min}}$ qui ont les propriétés suivantes

a) c'est une configuration, i.e. \inserted{\unsure{formé|formée} de fig. él.} \unsure{mutuellement|mutuellemt} \deleted{\unsure{dis}} compatibles, et elle est non vide

b) Elle est loc. finie.

Dans les exemples \unsure{types|type}, ceci correspond aux parties \inserted{fermées} $Y$ de $X$ qui sont des sous-variétés \inserted{non vides}, \unsure{découpées|découpées,} en strates par les composantes connexes.

Aussi, \unsure{\deleted{les}|les} figures de dim $\leq 0$ s'appeleront les \uline{sous}-\uline{variétés} (au sens strict, i.e. avec stratification triviale).

La sous-variété vide \unsure{correspond} \unsure{à|:} la configuration vide.
\scanpage{14}
\mspage{13}
\uline{Subdivision} d'une figure. \deleted{\unsure{Et|S}} On se reporte au cas-type. Soit une figure $F$, \unsure{formée} de $X$ et de \deleted{\unsure{l'ens|l'ens.}} de ses strates $X_i$.

Une \deleted{subdivision} figure $F'$ qui est une subdivision de $\deleted{F}\unsure{X|F}$ est \deleted{\unsure{une}} définie par les propriétés

\[
\text{a)}\quad |F'|=|F|
\]

b) la relation d'équiv. dans $X=|F|$ définie par $F'$ est plus fine que celle définie par \deleted{\unsure{$F$}} $F$.

La condition b) implique

b$'$) $\forall X_i$, strate (\unsure{fermée}) de $X$ pour $F$, $X_i$ est réunion de strates de $X$ pour $F'$, \deleted{\unsure{e|et}} et \inserted{elle} lui est \unsure{$=$|$\overline{\mathrm{m}}$} \unsure{équivalente}.

Ainsi, les strates pour $F'$ qui sont $\subset X_i$, forment une \inserted{figure, qui est une} subdivision de $F_{X_i}$ (la figure \inserted{élémentaire} induite par $F$ sur $X_i$)

Mais on fera attention que cette figure, de support $X_i$, n'est pas \unsure{néc.|néc} élémentaire (une subdivision d'une figure élém. peut n'être pas élémentaire)
\scanpage{15}
\mspage{14}

On va dire les choses autrement, en \unsure{introduisant} une autre relation d'ordre dans $\unsure{\mathcal{F}|F}_{el}$, plus fine que l'inclusion.

\deleted{\unsure{Soit|So}}

Soient $F_{\deleted{\unsure{0|e}}}^{\deleted{\unsure{F|T|\prime}}},\deleted{\unsure{F|E}}F_{\deleted{0}}'$ deux fig. \deleted{élémentaires}. \deleted{\unsure{On|Je}}Je dis que $F_{\deleted{0}}'\ll F_{\deleted{0}}^{\deleted{\prime}}$ (que $F_{\deleted{0}}'$ raffine $F_{\deleted{0}}^{\unsure{\deleted{\prime}|\deleted{T}}}$) si \unsure{\deleted{toute}|\deleted{tte}|toute|tte} \deleted{\unsure{cellule|celle}} strate \uline{ouverte} de $F_{\deleted{0}}'$ est contenue dans une strate ouverte de $F_{\deleted{0}}$ \inserted{(unique, bien sûr)}.

(\uline{NB} On aura \deleted{\unsure{supp.}} donc $|F'|\subset|F|$, mais pas \unsure{néc.|nec.} $|F'|=|F|$)

C'est une relation \uline{d'ordre} sur l'ens des figures, plus forte que celle d'inclusion (qui \deleted{\unsure{dit|est}} \inserted{signifie} que \unsure{\deleted{toute}|\deleted{tte}|toute|tte} \deleted{cellule}\inserted{strate} ouverte de $F'$ \uline{est} une \deleted{cellule}\inserted{strate} ouverte de $F$). Cette relation de \unsure{raffin|raffin.|rgle} \deleted{peut} \deleted{doit} se \unsure{déduire} \unsure{formellement|formellt} de la relation d'ordre \unsure{de|\deleted{de}} décrite sur $\mathcal{F}_{el}$, l'ens des figures élémentaires, car $F'$ raffine $F$ ssi pour tte

\[
\begin{aligned}
&\deleted{\unsure{F'_{X'}|F'_{X}}}\ (X'\in{\unsure{\Sigma|\overline{Z}}}^{\unsure{\prime|\deleted{\prime}}}_{F'}),\quad \exists X\in\Sigma^{\unsure{\prime|\deleted{\prime}}}_F,\ \text{telle que}\\
&\qquad F'_{X'}\ \text{raffine}\ F_X
\end{aligned}
\]
\scanpage{16}
\mspage{15}

Ceci dit, une figure $F'$ est une \uuline{subdivision} d'une figure $F$ ssi elle \uline{raffine} $F'$, et \inserted{si} \uline{elle a \unsure{$\overline{\mathrm{m}}$|mê|$=$} support}. \inserted{Mais (cf plus bas) cela \unsure{équivaut} \unsure{:|à} dire que $F$ est \deleted{\unsure{maximal|maximale}} \inserted{maximale} pour la relation $\leq$, \unsure{\deleted{parmi}|parmi} \deleted{\unsure{toutes les|tous les}} les figures qui raffinent $F$.}

Ainsi, on trouve d'autres éléments de structure :

\textcircled{5} Une relation d'ordre $\ll$ sur $\unsure{\mathcal{F}|F}_{el}$ (``raffinement''), impliquée par $\leq$, d'où une notion, pour deux figures $F\subset\mathcal{F}_{el}$ et $F'\subset\mathcal{F}_{el}$, de $F$ raffine $F'$ :

\[
\begin{aligned}
&\deleted{\unsure{\exists|F}}\ \forall\,\deleted{\unsure{F|X}}X\in\deleted{\unsure{\mathcal{F}_{el}|F_{el}}}F,\quad \exists X'\in F'\quad\text{telle que}\\
&X\ll X'.
\end{aligned}
\]

Ceci posé

\inserted{\uline{Fig}$_{\deleted{\unsure{2}}3}$ Pour $X,X'\in\mathcal{F}_{el}$, on : $X\leq X'$ ssi $\{X,X'\}$ comp. et $X\ll X'$.}

% slanted
\margin{Donc les structures \textcircled{3} et \textcircled{5} \unsure{impliquent} la structure \textcircled{1}}

\deleted{\textcircled{6} \unsure{1|I} : On appellera \unsure{de|la} subdivision d'une figure}

\inserted{\uline{\unsure{Conséquence}} Pour deux \unsure{\deleted{figures $F,F'$}|figures $F,F'$}$\subset\mathcal{F}_{el}$, on a $F\unsure{\subseteq|\leq}F'\unsure{\Longleftrightarrow|\Longrightarrow} F,F'$ compatibles, et $F\ll F'$}

Ceci posé, on dira qu'une figure $F'$ est une \uline{subdivision} \unsure{de' une|d'une} autre $F$, si c'est un \uline{raffinement}, et si elle est maximale, pour la relation $\subset$, parmi les raffinements de $F$.
\scanpage{17}
\mspage{16}

\uline{Justification} (?)

% slanted
\margin{\unsure{À|A} développer -- ou bien, il est \unsure{douteux} que la notion de raffinement \deleted{\unsure{im|in}}\unsure{implique} \unsure{déjà|déjà,} celle de subdivision. On va pour le moment l'admettre \unsure{formel|provi}\deleted{\unsure{lement|soirement}}. C'est beaucoup plus trivial que ça -- s'il existait une strate ouverte $U_i$ de $F$ telle que $U_i\not\subset X'$, il suffirait de prendre \deleted{la figure} $x\in U_i\setminus X'\cap U_i$, et la figure $F''=F'\cup\{x\}$ qui raffine $F$, et telle que $F'<F''$ !!}

Revenons au cas type, soit $F'$ un raffinement \inserted{de $F$}, tel que $|F'|\neq|F|$. Montrons que $F'$ n'est \unsure{pas} maximale, \inserted{pour $\leq$,} parmi les raffinements de $F$.

On va définir un raffinement \unsure{canonique} \inserted{$\overline{F'}$} de $F$, qui est \unsure{$\underline{\mathrm{m}}$|$\overline{\mathrm{m}}$|$\subseteq$|$=$} une subdivision, et qui contient la figure $F'$. (donc strictement, puisque $\underbrace{|F'|}_{X'}\neq\underbrace{|F|}_{X}=|\overline{F'}|$\unsure{.)|)} Les \deleted{\unsure{cellules|cellule}} strates ouvertes de $F''$ seront

a) les strates ouvertes \inserted{$U'_j$} de $F'$

b) Les comp. connexes \inserted{$V_{i,\alpha}$} des

\[
\begin{aligned}
&\deleted{\unsure{V_{i,\alpha}|V_{ij}}}\ \underbrace{U_i\setminus U_i\cap X'}\qquad \text{(\unsure{où|ou} } U_i \text{ est une}\\
&\hspace{14em}\text{strate ouverte de } \deleted{\unsure{X|E}}F\text{)}
\end{aligned}
\]

(i.e. de \unsure{:|.} ceux $\uparrow$ qui ne sont pas $\varnothing$)

\deleted{\unsure{Les|les}} Il faut prouver les conditions d'une stratification, et pour commencer que $\partial V_{i\alpha}\overset{\mathrm{def}}{=}\overline{V_{i\alpha}}\setminus V_{i,\alpha}$ est réunion de strates

\[
\text{Donc déjà pour}\quad \partial V_{i\alpha}\cap U_i\ \left\{\begin{aligned}&=\text{\deleted{\unsure{B|E}}Frontière de } V_{i\unsure{2|\alpha}}\text{ dans }U_i\text{.}\\[-2pt]&\subset X'\cap U_i\end{aligned}\right.
\]

il y a un pb, qui a l'air \unsure{intéressant|intimement}, de topologie d'un espace stratifié $(X'\cap U_i)$ \unsure{plongé} dans une va\uline{riété} (top. \unsure{ou} \unsure{modérée}) $(U_i)$
